\uuid{bigC}
\titre{Module/argument}
\auteur{bodin ; hanani ; mzari}
\organisation{exo7}
\licence{CC BY-NC-SA 4.0}
\datecreate{2020-01-20}
\niveau{L1}
\module{Algèbre}
\chapitre{Nombres complexes}
\sousChapitre{Forme cartésienne, forme polaire}
\difficulte{2}
\theme{Écritures algébrique et géométrique, module, argument}
\exo7qcmid{lille-1.nombres-complexes.007}

\contenu{
\begin{question}
\qclassification{104.01}
\qtags{motcle=module/argument}
Soit $z=\frac{(1-i)^{10}}{(1+i\sqrt 3)^4}$. Quelles sont les assertions vraies ?
\begin{answers}
    \good{$|z|=2$}

    \bad{$|z|=\frac{1}{2}$}

    \good{$\arg z = \frac{\pi}{6} \, [2\pi]$}

    \bad{$\arg z = -\frac{\pi}{6} \, [2\pi]$}

\end{answers}
\begin{explanations}
On applique les formules :
$|\frac{z_1^n}{z_2^m}|= \frac{|z_1|^n}{|z_2|^m}$   et $\arg(\frac{z_1^n}{z_2^m})= n\arg z_1 - m\arg z_2 \, [2\pi]$.
\end{explanations}
\end{question}
}
